\documentclass[11pt]{book}

\usepackage{notes}
\usepackage{math}

\begin{document}

% \var 
% \norm{}
% Nc \mathcal{N}
% \bp{}

\newcommand{\Cov}{\mathrm{Cov}}

\theoremstyle{definition}
\newtheorem{exmp}{Example}[section]

\chapter{February 10, 2026 Lecture}
\section{Maximum-Likelihood Estimation}
Finding an UMVUE can be mathematically challenging, and in some cases, it might not even exist. More importantly, while an UMVUE minimizes variance 
among unbiased estimators, strictly adhering to the unbiased constraint can somestimes lead to a significantly higher variance than a slightly biased estimator.
\subsubsection{Illuminating examples}
Perhaps the best known example 
\[
  X \sim Poisson(\lambda),
\]
where $\lambda > 0$ is the unknown parameter. If our goal is to estimate 
\[
  g(\theta) = \exp(-2\lambda),
\] 
then it can be shown that the only unbiased estimator is give by the laughable 
\[
  \hat{g}(X) = (-1)^X 
\]
Another illuminating example is to estimate the location of a light souce based on its projection on the boundary of 
a dart board over a randomly placed dart.
\section{Maximum-likelihood estimator (MLE)}
An alternative frequentist estimator is the maximum-likelihood estimator (MLE). Focusing on estimating the model parameter $\theta$ itself. Given a likelihood 
family $f_\theta(x)$, an estimator $\hat{\theta} = \hat{\theta}(X)$ is an MLE of $\theta$ if 
\[
  \hat{\theta}(x) \in \arg\max_{\theta\in\Theta} f_\theta(x), \quad \forall x\in\Xc
\]
or, equivalently,
\[
  \hat{\theta}(x) \in \arg\max_{\theta\in\Theta} \log f_\theta(x), \quad \forall x\in\Xc
\]
\subsection{Maximizing likelihood and minimizing K-L divergence}
While mainly motivated by intuition, MLE has an interesting interpretation when the observation $X = (X_1, \dots, X_n)$, where 
\[
  X_i \iid p_\theta(x_i)
\]
More specifically, given an observation $x = (x_1, \dots, x_n)$, let 
\[
  \hat{p}_x (z) = \frac{1}{n}\sum_{i=1}^n 1_{\{x_i=z\}}, \quad \forall z \in\Xc 
\]
be the \textit{empirical} estimate of the \textit{marginal} distribution given the observation $x$. For two pmf's $p$ and $q$ 
over the same alphabet $\Xc$, the \textbf{K-L divergence} between $p$ and $q$ is defined as: 
\[
  D_{KL}(p || q) := \E_p \left[ \log\frac{p(Z)}{q(Z)} \right] = \sum_{z\in\Xc} p(z) \log \frac{p(z)}{q(z)}
\]
\subsubsection{Information inequality}
Occurs when 
\[
  D_{KL}(p || q) \geq 0 
\]
where the equality holds if and only if $p=q$. In statistics, information theory, and machine learning, one often use K-L divergence as a distance measure between two probability
distributions, even though this distance measure is generally asymmetric about the two input distributions. \\ 

\noindent Claim: maximizing the likelihood 
\[
  \sum_{i=1}^n \log p_\theta(x_i)
\] 
is equivalent to minimzing the K-L divergence 
\[
  D_{KL}(\hat{p}_x || p_\theta)
\]
\begin{proof}
  First note that 
  \begin{align*}
    D_{KL}(\hat{p}_x || p_\theta) &= \E_{\hat{p}_x} \left[ \log \frac{\hat{p}_x(Z)}{p_\theta(Z)}\right] \\ 
                                  &= \E_{\hat{p}_x} [\log\hat{p}_x(Z)] - \E_{\hat{p}_x}[\log p_\theta(Z)]
  \end{align*}
  Further note that 
  \begin{align*}
  \E_{\hat{p}_x} &= \sum_{z\in\Xc} \hat{p}_x(z)\log p_{\theta}(z) \\ 
                   &= \frac{1}{n} \sum_{z\in\Xc}\sum_{i=1}^n 1_{\{ x_i=z \}} \log p_\theta(z) \\ 
                   &= \frac{1}{n} \sum_{i=1}^n \bp{\sum_{z\in\Xc} 1_{\{x_i=z\}} \log p_\theta(z)} \\ 
                   &= \frac{1}{n} \sum_{i=1}^n \log p_\theta(x_i)
  \end{align*}
  Combining the above two results gives 
  \[
    D_{KL}(\hat{p}_x || p_\theta) = \E_{\hat{p}_x}[\log \hat{p}_x (Z)] - \frac{1}{n} \sum_{i=1}^n \log p_\theta(x_i)
  \]
  We may thus conclude that minimizing the K-L divergence $D_{KL}(\hat{p}_x || p_\theta)$ is equivalent to 
  maximizing the likelihood $\sum_{i=1}^n \log p_\theta(x_i)$. The above result also suggests that if 
  our model is identifiable, MLE will recover $\theta$ in the limit as $n \to \infty$. We will have more 
  to say about this result a bit later.
\end{proof}

\subsection{Re-parameterization}
Assume that the likelihood function $f_\theta(x)$ can be re-parameterized in terms of a more natural parameter 
\[
  \eta = \eta(\theta)
\]
In this case, for any given observation $x$, we have 
\[
  \arg\max_{\theta\in\Theta} f_\theta(x) = \{ \theta \in \Theta : \eta(\theta) \in \arg\max_{\theta\in\Theta} \}
\] ... %TODO
Consider an exponential family of distribution re-parameterized by the natural parameter $\eta$: 
\[
  \tilde{f}_\eta(x) = h(x) \exp(\eta^t T(x) - A(\eta))
\]
where the space of the natural parameter is an open and connected subset of $\R^s$. To find an MLE for $\eta$, note that 
\begin{align*}
  \nabla_\eta \log\tilde{f}_\eta(x) &= T(x) - \nabla_\eta A(\eta) = T(x) - \E_\eta[T(X)] \\ 
  \nabla_\eta^2 \log \tilde{f}_\eta(x) &= -\nabla_\eta^2 A(\eta) = \Cov_\eta[T(X)]
\end{align*}
A covariance mastrix is always positive semi-definite, so for any $x\in\Xc$, the log-likelihood $\log\tilde{f}_\eta(x)$ is a concave function of $\eta$, for which any stationary point is a global maxima. Therefore,
any solution for $\eta$ to the equation 
\[
  \E_\eta[T(X)] = T(x)
\] 
is an MLE for $\eta$. Furthermore, if the covariance matrix of $T(X)$ is positive definite for all 
$\eta$, for any $x \in \Xc$ the log-likelihood $\log\tilde{f}_\eta(x)$ is a strictly concave function of $\eta$. In this case, the MLE is unique if it exists.

Finally, by the re-parameterization property, any solution for $\theta$ to the equation 
\[
  \E_\theta[T(X)] = T(x)
\]
is an MLE for $\theta$.
\subsubsection{Gaussian with unknown mean and variance}
Recall that in this case, a sufficient statistic is given by 
\[
  T_1 = \sum_{i=1}^n X_i \quad\text{and}\quad T_2 = \sum_{i=1}^n X_i^2
\]
the expectations of the sufficient statistics are given by: 
\[
  \E_\theta[T_1] = n\mu \quad\text{and}\quad \E_\theta[T_2] = n(\sigma^2 + \nu^2)
\]
From previous discussion, any solution of $(\mu, \sigma^2)$ to the equations 
\begin{align*}
  n\mu &= T_1(x) \\ 
  n(\sigma^2 + \mu^2) &= T_2(x)
\end{align*}
is an MLE. The above equations have a unique solution given by 
\[
  \mu = \frac{T_1(x)}{n} \quad\text{and}\quad \sigma^2 = \frac{T_2(x)}{n} - \bp{\frac{T_1(x)}{n}}^2
\]
  We may thus conclue that the empirical mean 
  \[
    \hat{\mu} = \frac{T_1(x)}{n} = \frac{1}{n} \sum_{i=1}^n X_i
  \]
  is the MLE for $\mu$ and the empirical variance 
  \[
    \hat{\sigma}^2 = \frac{T_2(x)}{n} - \bp{\frac{T_1(x)}{n}}^2 = \frac{1}{n}\sum_{i=1}^n X_i^2 - \hat{\mu}^2 = \frac{1}{n}\sum_{i=1}^n (X_i - \hat{\mu})^2 
  \]
  is the MLE for $\sigma^2$, when both $\mu$ and $\sigma^2$ are unknown. By Cochran's theorem,
  \[
    \E_\theta \left[ \sum_{i=1}^n (X_i - \hat{\mu})^2\right] = (n-1)\sigma^2 
  \]
  so 
  \[
    \E_\theta[\hat{\sigma^2}] = \frac{n-1}{n}\sigma^2 
  \] 
  i.e., the empirical variance is a biased estimator for the true variance $\sigma^2$.

\subsection{Consistency and asymptotic normality}
Note that in the previous example, while the MLE for $\sigma^2$ is biased, the bias vanishes asymptotically in the limit as the sample size $n\to\infty$. Also note that the variances of the MLE for both $\mu$ and $\sigma^2$ coincide with the CRLB for unbiased estimators. Under certain regularity conditions, it can be shown that the above observations hold for a general likelihood function.

More specifically, let $X = (X_1, \dots, X_n)$, where $X_i \iid f_\theta(x_i)$. For any $\theta$ that is in the interior of $\Theta$, let $\hat{\theta}$ be the MLE for $\theta$ given by the unique solution to the equation 
\[
  \nabla_\theta \bp{\sum_{i=1}^n \log f_\theta(x_i)} = 0
\] 
We have $\sqrt{n}(\hat{\theta} - \theta)$ converges in distribution to $\Nc(0, I_1^{-1}(\theta))$ in the limit as the sample size $n\to\infty$.
That is, when the sample size $n$ is large, the MLE $\hat{\theta}$ is approximately distributed as 
\[
  \Nc \bp{\theta, \frac{1}{n}I_1^{-1}(\theta)}
\] 
or, equivalently,
\[
  \Nc (\theta,I_n^{-1}(\theta))
\] 
due to the additivity of Fisher information $I_n(\theta) = n I_1(\theta)$. Therefore, under certain regularity conditions, the MLE is consistent in that it can recover the unknown parameter $\theta$ in the limit as the sample size $n\to\infty$. The proof of the asymptotic normality of the MLE is based on the well-known central limit theorem. This is omitted here, however.

\subsection{Multinomial with unknown probabilities}
Goal is to use the empirical counts to estimate the true probability. 
Let $X\sim p_\theta(x)$, where the observation 
\[
   X = (N_1, \dots, N_K)
 \]
 is the empirical counts over $K$ categories and $n$ total counts, the unknown parameter 
 \[
   \theta = (p_1, \dots, p_K)
 \] 
 is a probability vector over $K$ categories, and the likelihood function is multinomial: 
 \[
   p_\theta(x) = \frac{n!}{n_1!\cdots n_K!} \prod_{k=1}^K{p_k^{n_k}}
  \]
  % TODO: finish this whole section 
  The total counts $n$ and the number of categories $K$ are assumed to be known. The MLE can be found by 
  solving the equation 
  \[
    \nabla_\theta L_\lambda (\theta) = 0 
  \] 
  and then finding a Lagrangian multiplier $\lambda \geq 0$ to satisfy the constraint 
  \[ 
    \sum_{k=1}^K p_k = 1 
  \]
  The MLE for $\theta$ in this case is given by the empirical frequency estimator 
  \[
    \hat{\theta} = \bp{\frac{N_1}{n}, \dots, \frac{N_K}{n}}
  \] 
Note that for any $k \in [K]$ 
\[
  \E_\theta[N_k] = np_k 
\] 
so the empirical frequency estimator $\hat{\theta}$ is unbiased. 
\section{Estimating $g(\theta)$}
Let $\eta = g(\theta)$. Note that if the likelihood function $f_\theta(x)$ can be rep-parameterized in terms of 
$\eta$, then it is natural to use an MLE of $\eta$ as an "MLE" estimate of $g(\theta)$. By the re-parameterization 
property, $\eta^*$ is an MLE of $\eta$ if and only if 
\[
  \eta^* = g(\theta^*)
\] 
for some $\theta^*$ that is an MLE of $\theta$. Therefore, if the likelihood function $f_\theta(x)$ can be re-parameterized 
in terms of $g(\theta)$, the "MLE" of $g(\theta)$ is given by $g(\theta^*)$, where $\theta^*$ is an MLE of $\theta$. 

In statistics, it is customary to define the MLE of $g(\theta)$ as $g(\theta^*)$, whether the likelihood function $f_\theta(x)$ can be re-parameterized
in terms of $g(\theta)$ or not. This is known as the invariance property of MLE. Many other estimators 
such as UMVUE do not enjoy the invariance property. 
%\chapter{February 12, 2026 Lecture}
   %
\chapter{The Expectation-Maximization Algorithm}
In general, finding MLE requires iterative numerical techniques. Gradient descent methods for 
minimization as Newton-Raphson and Fisher's scoring methods are commonly used. Here we discuss a method known as 
expectation-maximization (EM) algorithm, which is specifically designed for computing MLE for models with latent variables. 
\section{Models with latent variables}
Let $(X,U) \sim f_\theta(x, u)$, where $X$ is the observation and $U$ is the unobserved/latent variable. We seek the ML estimate of $\theta$ based on the observation $X$: 
\[
  \hat{\theta}(x) \in \arg\max_{\theta\in\Theta} \log f_\theta(x)
\]
where 
\[
  f_\theta(x) = \int_{\Uc} f_\theta(x, u) \, du 
\] 
The above obtimization problem can be difficult to solve analytically due to the marginalization of the latent variables, 
which occurs inside the logarithm of the objective function.
\subsection{Mixture Gaussian model}
Let $X = (X_1, \dots, X_n)$ and $U = (U_1, \dots, U_n)$, where $(X_i, U_i)$ are i.i.d. according to: 
\[
  U_i \sim Cat(u_i; \pi_1, \dots, \pi_K) \quad\text{and}\quad X_i|U_i \sim \Nc(x_i; \mu_{i_i}, \sigma_{u_i}^2)
\]
Marginalizing over $U_i$ gives 
\[
  f_\theta(x_i) = \sum_{k=1}^K \pi_k \Nc(x_i; \mu_k, \sigma_k^2)
\] 
which is a mixture Gaussian distribution, and 
\[
  \theta = (\pi_1, \dots, \pi_k, \mu_1, \dots, \mu_K, \sigma_1^2, \dots, \sigma_K^2) 
\] 
is the unknown parameter. The log-likelihood for all measurements is given by 
\[
  \log f_\theta(x) = \sum_{i=1}^n \log f_\theta(x_i) = \sum_{i=1}^n \log \bp{\sum_{k=1}^K \pi_k \Nc (x_i; \mu_k, \sigma_k^2)}
\]
\subsubsection{Finding MLE for $(\mu_1, \dots, \mu_K)$}
Assume for the moment that $(\pi_1, \dots, \pi_K)$ and $(\sigma_1^2, \dots, \sigma_K^2)$ are known. To find the MLE 
for $(\mu_1, \dots \mu_K)$, note that for any $j \in [K]$ we have 
\[
  \frac{\partial}{\partial \mu_j} \log f_\theta(x) = \sum_{i=1}^n \frac{\tilde{\pi}_j(x_i)(x_i - \mu_j)}{\sigma_j^2}
\]
where 
\[ 
  \tilde{\pi}_j (x_i) = \frac{x_i \Nc(x_i; \mu_j, \sigma_j^2}{\sum_{k=1}^K x_i \Nc (x_i; \mu_k, \sigma_k^2)} = \mathbb{P}_\theta(U_i = j | X_i = x_i)
\]
Setting the partial derivatives equal to zero for all $j$ gives the following set of equations: 
\[
  \sum_{i=1}^n \frac{\tilde{\pi_j}(x_i)(x_i - \mu_j)}{\sigma_j^2} = 0, \quad\forall j\in[K]
\]
for which the solution for $(\mu_1, \dots, \mu_K)$ is the MLE.
\subsubsection{A vicious circle}
Note that if $\tilde{\pi}_j(x_i)$'s are known, the above equations are separate for each 
$\mu_j$ and hence are very easy to solve: % TODO: here, replace pi with x_i (he is updating notes )
\[
  \mu_j = \frac{\sum_{i=1}^n \tilde{\pi}_j(x_i)x_i}{\sum_{i=1}^n \tilde{\pi}_j(x_i)}, \quad\forall j\in[K]
\]
However, in reality, $\tilde{\pi}_j(x_i)$'s depend on $(\mu_1, \dots, \mu_K)$ and hence are unknown. Therefore,
the above equations are interlocked and hence are very difficult to solve. This is a common challenge for 
finding the MLE for models with latent variables. We are now in the following situation:
\begin{enumerate}
  \item If we know $\tilde{\pi}_j(x_i)$'s, we can solve for $\mu_j$'s 
  \item If we know $\mu_j$'s, we can compute $\tilde{\pi}_j(x_i)$'s 
\end{enumerate}
\section{An iterative procedure}
A successive approximation procedure: 
\begin{enumerate}  
  \item Initialize $(\mu_1, \dots, \mu_K)$
  \item Evaluate the posterior probabilities $(\tilde{\pi}_1(x_i), \dots, \tilde{\pi}_K(x_i)), i \in [n]$ using the current value of $(\mu_1, \dots, \mu_K)$
  \item Update the value for $(\mu_1, \dots, \mu_K)$ using the posterior probabilities from Step 2 
  \item Stop if the value for $(\mu_1, \dots, \mu_K)$ has already converges; otherwise go back to Step 2 
\end{enumerate}
Next, we shall generalize the above prodedure to general likelihood models with latent variables, and the resulting procedure is known as 
the expectation-maximization (EM) algorithm. 
\section{A variational characterization of the log-likelihood}
The main idea behind the EM algorithm is the following variational characterization of the log-likelihood:
\[
  \log f_\theta(x) = \max_{q} F(q, \theta)
\] 
where 
\begin{align}
  F(q,\theta) &= \log f_\theta(x) - D_{KL}(q(u)|| f_\theta(u|x)) \\ 
              &= \E_q[\log f_\theta(x, U)] - \E_q[\log q(U)],
\end{align}
and both expectations are over $U \sim q(u)$. The main advantage of this variational characterization
is that it allows us to rewrite the problem of finding the MLE as a double-maximization problem: 
\[ 
  (\hat{\theta}, q^*) \in \arg \max_{(\theta, q)} F(q, \theta) 
\]
\begin{proof}
  First note that by the information inequality, for any density function $q$ we have 
  \[
    D_{KL}(q(u) || f_\theta(u|x)) \geq 0 
  \]
  where the equality holds if and only if 
  \[
    q(u) = f_\theta(u | x), \quad\forall u \in\Uc 
  \]
  We thus have
  \[
    \log f_\theta(x) \geq \log f_\theta(x) - D_{KL}(q(u) || f_\theta(u | x))
  \] 
  where the equality holds if and only if 
  \[
    q(u) = f_\theta(u | x), \quad\forall u \in \Uc 
  \] 
  and hence 
  \[
    \log f_\theta(x) = \max_{q} (\log f_\theta(x) - D_{KL}(q(u) || f_\theta(u|x)))
  \] 
  Further note that 
  \begin{align*}
    \log f_\theta&(x) - D_{KL}(q(u) || f_\theta(u | x)) \\
                 &= \log f_\theta(x) - (\E_q[\log q(U)] - \E_q[\log f_\theta(U|x)]) \\ 
                 &= \E_q [\log f_\theta(x)] + \E_q[\log f_\theta(U|x)] - \E_q[\log q(U)] \\ 
                 &= \E_q [\log (f_\theta(x)f_\theta(U|x))] - \E_q[\log q(U)] \\ 
                 &= \E_q[\log f_\theta(x, U)] - \E_q[\log q(U)]
  \end{align*}
\end{proof}
We have thus proved that (2.1) and (2.2) are in fact equivalent to each other. 

\section{Finding MLE via alternating maximization}
Recall that the MLE can be found by solving the double-maximization problem: 
\[
  \max_{(\theta, q)} F(q, \theta) 
\]
Consider the following alternating-maximization strategy for solving the above double-maximization problem: 
\begin{enumerate}
  \item Initialize $\theta = \theta^{(0)}$ and $t=1$
  \item Repeat the following steps until the values of $(\theta, q)$ converge:
    \begin{enumerate}
      \item Fix $\theta = \theta^{(t-1)}$ and find 
        \[
          q^{(t)} \in \arg\max_{q}F\bp{q,\theta^{(t-1)}}
        \] 
      \item Fix $q = q^{(t)}$ and find
        \[
          \theta^{(t)} = \arg\max_{\theta}F\bp{q^{(t)}, \theta}
        \]
      \item $t \leftarrow t + 1$
    \end{enumerate}
\end{enumerate}
The problem of fixing $\theta = \theta^{(t-1)}$ and finding 
\[
  q^{(t)} \in \arg\max_{q}F\bp{q,\theta^{(t-1)}}
\] 
is trivial. By the information inequality, 
\[
  q^{(t)}(u) = f_{\theta^{(t-1)}} (u|x),\quad\forall u\in\Uc 
\] 
which gives 
\[
  F\bp{q^{(t)}, \theta^{(t-1)}} = \log f_{\theta^{(t-1)}}(x)
\]
By comparison, the problem of fixing $q=q^{(t)}$ and finding 
\[
  \theta^{(t)} \in \arg\max_{\theta}F\bp{q^{(t)},\theta} 
\]  
is much more involved. To solve this optimization problem, we shall use expression (2) and write % TODO: fix eqn reference here and above 
\[
  F\bp{q^{(t)}, \theta} = \E_{q^{(t)}} [ \log f_\theta(x, U)] - \E_{q^{(t)}} [\log q^{(t)}(U)]
\]
Note that the second term in the above expression does \textit{not} involve $\theta$. Thus, $\theta^{(t)}$ can be found by maximizing
over the first term in the above expression 
\[
  \theta^{(t)} \in \arg\max_{\theta} \E_{q^{(t)}} [\log f_\theta(x, U)] 
\] 
where 
\[
  q^{(t)} (u) = f_{\theta^{(t-1)}} (u | x), \quad\forall u\in\Uc 
\]
The above two equations allow us to continuously update the value of $(\theta, q)$, leading to the famous EM algorithm.




\end{document}
